% Overview figure for the Times New Roman revision: the same diagram as overview_tikz.tex, with the three math items
% (lambda, k, G) written as italic Unicode letters so that every glyph comes from Times New Roman.
% Compile with XeLaTeX (overview_times_standalone.tex).
% Overview figure: probe-based analysis of adversarial feature removal.
% Preamble requirements:
%   \usepackage{tikz}
%   \usetikzlibrary{arrows.meta,calc}
\begin{tikzpicture}[
  font=\footnotesize,
  >={Stealth[length=2.2mm]},
  box/.style={draw=black!80, line width=0.5pt, rounded corners=2pt,
              align=center, inner sep=4pt, minimum height=12mm,
              text width=#1-8pt, minimum width=#1,
              execute at begin node={\hyphenpenalty=10000\relax}},
  box/.default=4cm,
  std/.style={dashed, fill=white},
  ours/.style={fill=black!8},
  fail/.style={fill=black!22, line width=0.9pt},
  arr/.style={->, line width=0.6pt, draw=black!85},
]
% Column grid: total width 15 cm
\def\W{15}
\def\half{7.25}   % half-width boxes (0.5 cm gap)
\def\third{4.6667} % third-width boxes (0.5 cm gaps)

% Row 1: data and preprocessing
\node[box=\half cm, anchor=north west] (data) at (0,0)
  {\textbf{Four public smartphone IMU datasets}\\
   HHAR, RealWorld, MotionSense, MobiAct};
\node[box=\half cm, anchor=north east] (prep) at (\W,0)
  {\textbf{Leak-free preprocessing}\\
   50\,Hz resampling, 2.56\,s windows,\\segment-level splits};
\draw[arr] (data) -- (prep);

% Row 2: adversarial training (container drawn explicitly)
\def\ya{-2.2}  \def\yb{-5.15}
\draw[black!60, dashed, rounded corners=4pt, line width=0.5pt]
  (0,\ya) rectangle (\W,\yb);
\node[anchor=north west, font=\footnotesize\bfseries, inner sep=0pt]
  at (0.3,\ya-0.25)
  {Adversarial training over a \textit{λ} grid, all seeds and conditions};
\draw[arr] (prep.south) -- (prep.south |- 0,\ya);

\node[box=3.6cm, anchor=north west] (enc) at (0.4,\ya-0.8)
  {\textbf{1D CNN encoder}\\TCN, Transformer and BiLSTM as replications};
\node[box=4.4cm, anchor=north] (tgt) at (7.5,\ya-0.8)
  {\textbf{Target head}\\identity\\(activity in the swap experiment)};
\node[box=4.4cm, anchor=north east] (disc) at (\W-0.4,\ya-0.8)
  {\textbf{Discriminator}\\behind gradient reversal\\activity, device, position\\(subject in the swap experiment)};
\draw[arr] (enc.east |- 0,-3.4) -- (tgt.west |- 0,-3.4);
\coordinate (clearrow) at ($(tgt.south)+(0,-0.15)$);
\draw[arr] (enc.south) -- (enc.south |- clearrow) -- (disc.west |- clearrow);

% Row 3: two ways to judge invariance
\def\yc{-6.0}
\node[box=6.9cm, std, anchor=north west] (stdchk) at (0,\yc)
  {\textbf{Training-time check}\\discriminator accuracy falls to chance};
\node[box=6.9cm, ours, anchor=north east] (probe) at (\W,\yc)
  {\textbf{Post-hoc probe check}\\
   linear, MLP, \textit{k}-NN fitted to the frozen encoder,\\
   scored on held-out segments, devices and subjects};
\draw[<->, line width=0.6pt, draw=black!85] (6.9,\yc-0.6) -- node[above, font=\scriptsize, inner sep=1pt] {\textit{G}} (\W-6.9,\yc-0.6);
\draw[arr] (stdchk.north |- 0,\yb) -- (stdchk.north);
\draw[arr] (probe.north |- 0,\yb) -- (probe.north);

% Row 4: findings
\def\yd{-8.4}
\node[box=\third cm, fail, anchor=north west] (f1) at (0,\yd)
  {\textbf{Failure 1: discriminator reads low}\\
   fails to keep up with the encoder; collapse is the extreme case, reading chance while the factor stays decodable};
\node[box=\third cm, fail, anchor=north] (f2) at (\W/2,\yd)
  {\textbf{Failure 2: linear gap}\\
   linear probes report far more removal than nonlinear probes};
\node[box=\third cm, ours, anchor=north east] (f3) at (\W,\yd)
  {\textbf{Secondary findings}\\
   training-set size, narrow band of \textit{λ}, weak link to decodability};

\coordinate (bus) at (0,-7.8);
\coordinate (fa) at ($(f1.north)+(-0.7,0)$);
\coordinate (fb) at ($(f1.north)+(0.7,0)$);
\draw[arr] (fa |- stdchk.south) -- (fa);
\draw[line width=0.6pt, black!85] (probe.south) -- (probe.south |- bus);
\draw[line width=0.6pt, black!85] (fb |- bus) -- (f3.north |- bus);
\draw[arr] (fb |- bus) -- (fb);
\draw[arr] (f2.north |- bus) -- (f2.north);
\draw[arr] (f3.north |- bus) -- (f3.north);

% Row 5: protocol
\def\ye{-11.6}
\node[box=\W cm, line width=0.9pt, anchor=north west] (proto) at (0,\ye)
  {\textbf{Reporting checklist}\\
   probe the frozen encoder with linear, MLP and \textit{k}-NN probes on held-out data;\\
   score the discriminator on held-out data (balanced accuracy); report headroom and the gap \textit{G};\\
   check for collapse (balanced accuracy, entropy, macro-F1); choose \textit{λ} on separate data;\\
   report the controls, the \textit{λ} sweep, the training-set size and the task cost};
\foreach \f in {f1,f2,f3} \draw[arr] (\f.south) -- (\f.south |- proto.north);

% Legend
\def\yl{-14.2}
\node[box=7mm, std, minimum height=4mm, inner sep=0pt, anchor=west] (l1) at (0,\yl) {};
\node[anchor=west, font=\scriptsize] (t1) at (l1.east) {Check available during training};
\node[box=7mm, ours, minimum height=4mm, inner sep=0pt, anchor=west] (l2) at ($(t1.east)+(0.5,0)$) {};
\node[anchor=west, font=\scriptsize] (t2) at (l2.east) {Probe-based measurement and findings};
\node[box=7mm, fail, minimum height=4mm, inner sep=0pt, anchor=west] (l3) at ($(t2.east)+(0.5,0)$) {};
\node[anchor=west, font=\scriptsize] at (l3.east) {Failures of the training-time check};
\end{tikzpicture}

